\paragraph{Low-Rank Neural Networks.}
Low-rank factorizations have been widely used for parameter-efficient deep learning~\cite{arora2018optimization,novikov2015tensorizing}. However, most work focuses on compression and efficiency rather than understanding the learning mechanisms. Our work shows that low-rank is not just about efficiency but enables fundamental learning properties like hierarchical construction.

\paragraph{Mean-Field Theory.}
Mean-field theory provides a framework for analyzing infinite-width neural networks~\cite{mei2018mean,chizat2018global,rotskoff2018neural}. We leverage this framework to understand how rank constraints affect training dynamics, showing that extensive-rank dynamics naturally give rise to hierarchical construction.

\paragraph{Hierarchical Learning.}
Hierarchical feature learning has been studied in various contexts~\cite{zeiler2014visualizing,olah2017feature}. Our work connects hierarchical learning to rank constraints, showing that low-rank enforces hierarchical structure through sequential activation.

\paragraph{Spectral Bias and Frequency Learning.}
Neural networks exhibit spectral bias toward low frequencies~\cite{rahaman2019spectral,xu2019training}. We extend this to show that in the extensive-rank regime, frequency bands are learned sequentially, creating hierarchical structure.

\paragraph{Edge-of-Stability.}
The edge-of-stability (EOS) phenomenon~\cite{cohen2021gradient} describes training at the boundary of stability. We show that EOS is not just a curiosity but a mechanism that enables hierarchical construction through instability-driven transitions.

\paragraph{Multi-Index Models.}
Multi-index models learn functions of the form $f(x) = \sum_j g_j(\langle v_j, x \rangle)$~\cite{benarous2021online}. We show that extensive-rank networks naturally learn multi-index structure through sequential direction activation.
